\section{Kiến thức nền tảng}
\label{sec:kien-thuc-nen}
 
Chương này thiết lập ngôn ngữ toán học dùng xuyên suốt báo cáo. Kim chỉ nam xuyên suốt của Geometric Machine Learning là khái niệm đối xứng: những thuộc tính của dữ liệu được bảo toàn dưới tác động của một họ phép biến đổi~\cite{bronstein2021geometric}. Để mô hình hoá ý tưởng này một cách chặt chẽ, nhóm trình bày theo trình tự từ trực giác đến hình thức: trước hết là hai khái niệm cốt lõi bất biến và đẳng biến; tiếp đó là công cụ đại số dùng để định nghĩa chính xác họ phép biến đổi hay còn gọi là lý thuyết nhóm, bao gồm tác động nhóm, biểu diễn nhóm và quỹ đạo; cuối cùng quay lại phát biểu hình thức của bất biến và đẳng biến dưới ngôn ngữ nhóm.

\subsection{Khái niệm về đối xứng, bất biến và tương biến}
\label{subsec:doi-xung-bat-bien}
 
Một cách trực quan, đối xứng của một đối tượng là phép biến đổi tác động lên đối tượng đó nhưng giữ nguyên những thuộc tính mà ta quan tâm. Chẳng hạn, nếu $x$ là một bức ảnh chứa con mèo và $T$ là phép tịnh tiến con mèo sang vị trí khác trong khung hình, thì bản chất đây là con mèo hoàn toàn độc lập với vị trí của nó. Khi xây dựng một mô hình $f$ (ví dụ bộ phân loại, bộ phân đoạn ảnh, hay mạng nơ-ron đồ thị), ta mong muốn $f$ tôn trọng các đối xứng này. Có hai cách tôn trọng cơ bản, tương ứng với hai loại bài toán: bất biến và tương biến / đẳng biến.
 
\begin{definition}[Phép đối xứng]
\label{def:doi-xung}
Cho đối tượng $x$ thuộc không gian dữ liệu $\mathcal{X}$ và một phép biến đổi $T:\mathcal{X}\to\mathcal{X}$. Ta nói $T$ là một phép đối xứng đối với một thuộc tính $P$ nếu $P\bigl(T(x)\bigr)=P(x)$ với mọi $x$. Tập hợp tất cả các phép đối xứng đối với một họ thuộc tính cho trước tạo thành một nhóm, cấu trúc sẽ được định nghĩa ở mục~\ref{subsec:group-theory}~\cite{bronstein2021geometric}.
\end{definition}
 
\subsubsection{Tính bất biến}
\label{ssub:invariance}
 
\begin{definition}[Bất biến]
\label{def:bat-bien}
Hàm $f:\mathcal{X}\to\mathcal{Y}$ được gọi là bất biến đối với phép biến đổi $T$ nếu~\cite{bronstein2021geometric}
\begin{equation}
\label{eq:invariance}
f\bigl(T(x)\bigr) = f(x), \qquad \forall x\in\mathcal{X}.
\end{equation}
Nói cách khác, đầu ra của mô hình không thay đổi khi đầu vào bị biến đổi bởi $T$.
\end{definition}
 
\begin{example}[Phân loại ảnh: bất biến với phép tịnh tiến]
\label{vd:cho}
Gọi $x$ là ảnh đầu vào và $T=\tau_v$ là phép tịnh tiến nội dung ảnh theo vector $v$. Bộ phân loại $f$ cần thoả $f(\tau_v(x))=f(x)$: dù con chó nằm ở góc nào trong khung hình, nhãn dự đoán vẫn là chó.
\end{example}

\begin{figure}[!htbp]
\centering
\begin{tikzpicture}[font=\small,
   frame/.style={draw=black!45, rounded corners=2pt,
                 minimum width=1.6cm, minimum height=1.6cm},
   fbox/.style={draw=gmlblue, fill=gmlblue, text=white, rounded corners=3pt,
                minimum width=0.9cm, minimum height=0.9cm, font=\bfseries\large},
   outlbl/.style={draw=green!55!black, fill=green!12, rounded corners=2pt,
                  inner sep=4pt, font=\bfseries}]
  \node[frame] (f1) at (0,2.05) {};
  \dogface{-0.32}{1.78}{0.56}
  \node[anchor=south] at (f1.north) {$x$};
  \node[frame] (f2) at (0,-0.05) {};
  \dogface{0.32}{0.22}{0.56}
  \node[anchor=north] at (f2.south) {$\tau_v(x)$};
  \draw[-{Latex[length=2mm]}, gmlorange, very thick] (f1.south) -- (f2.north);
  \node[gmlorange, anchor=west, font=\itshape] at (0.18,1.0) {$\tau_v$ (tịnh tiến)};
  \node[fbox] (f) at (3.4,1.0) {$f$};
  \draw[-{Latex[length=1.8mm]}, thick] (f1.east) .. controls (1.9,2.05) and (2.4,1.4) .. (f.north west);
  \draw[-{Latex[length=1.8mm]}, thick] (f2.east) .. controls (1.9,-0.05) and (2.4,0.6) .. (f.south west);
  \node[outlbl] (o) at (5.5,1.0) {Chó};
  \draw[-{Latex[length=1.8mm]}, thick] (f.east) -- (o.west);
\end{tikzpicture}
\caption[Bất biến với phép tịnh tiến]{Bất biến với phép tịnh tiến.\protect\\[2pt]
\textit{Ghi chú: Dù con chó nằm ở vị trí nào trong khung hình ($x$ hay bản tịnh tiến $\tau_v(x)$), bộ phân loại $f$ vẫn cho cùng một nhãn, tức $f(\tau_v(x))=f(x)$.}}
\label{fig:bat-bien-anh}
\end{figure}

\FloatBarrier

\begin{example}[Phân loại đồ thị: bất biến với phép hoán vị]
\label{vd:graph-inv}
Một đồ thị $n$ đỉnh thường được biểu diễn bằng ma trận kề $A\in\mathbb{R}^{n\times
n}$, trong đó cách đánh số đỉnh là hoàn toàn tuỳ ý. Nếu ta đánh số lại các đỉnh bằng một hoán vị, biểu diễn chuyển thành $PAP^{\!\top}$ với $P$ là ma trận hoán vị, nhưng vẫn là cùng một đồ thị. Do đó một mô hình dự đoán thuộc tính ở mức đồ thị, ví dụ phân tử có độc hay không, cần bất biến hoán vị:
\begin{equation}
\label{eq:graph-inv}
f\bigl(PAP^{\!\top}\bigr) = f(A), \qquad \forall P\in\mathcal{P}_n .
\end{equation}

\begin{figure}[!htbp]
\centering
\begin{tikzpicture}[font=\small,
   nd/.style={circle, draw=gmlblue, fill=gmlblue!12, minimum size=6.5mm,
              inner sep=0pt, font=\small},
   fbox/.style={draw=gmlblue, fill=gmlblue, text=white, rounded corners=3pt,
                minimum width=0.9cm, minimum height=0.9cm, font=\bfseries\large},
   outlbl/.style={draw=green!55!black, fill=green!12, rounded corners=2pt,
                  inner sep=4pt, font=\bfseries}]
  \begin{scope}
    \node[nd] (a1) at (0,1) {1}; \node[nd] (b1) at (1,1) {2};
    \node[nd] (c1) at (1,0) {3}; \node[nd] (d1) at (0,0) {4};
    \draw[gmlblue!70] (a1)--(b1) (b1)--(c1) (c1)--(d1) (d1)--(a1) (a1)--(c1);
    \node[anchor=east] at (-0.35,0.5) {$A$};
  \end{scope}
  \begin{scope}[shift={(0,-2.4)}]
    \node[nd] (a2) at (0,1) {3}; \node[nd] (b2) at (1,1) {4};
    \node[nd] (c2) at (1,0) {1}; \node[nd] (d2) at (0,0) {2};
    \draw[gmlblue!70] (a2)--(b2) (b2)--(c2) (c2)--(d2) (d2)--(a2) (a2)--(c2);
    \node[anchor=east] at (-0.35,0.5) {$PAP^{\!\top}$};
  \end{scope}
  \draw[-{Latex[length=2mm]}, gmlorange, thick]
        (1.35,0.5) .. controls (1.9,0.0) and (1.9,-1.4) .. (1.35,-1.9);
  \node[gmlorange, font=\itshape\footnotesize, anchor=west] at (2.15,-0.7) {$P$ (đánh số lại)};
  \node[fbox] (f) at (5.5,-0.7) {$f$};
  \draw[-{Latex[length=1.8mm]}, thick] (1.05,0.5)  -- (4.0,0.5)  .. controls (4.7,0.4)  .. (f.north west);
  \draw[-{Latex[length=1.8mm]}, thick] (1.05,-1.9) -- (4.0,-1.9) .. controls (4.7,-1.8) .. (f.south west);
  \node[outlbl] (o) at (7.4,-0.7) {``Độc''};
  \draw[-{Latex[length=1.8mm]}, thick] (f.east) -- (o.west);
\end{tikzpicture}
\caption[Bất biến với phép hoán vị]{Bất biến với phép hoán vị.\protect\\[2pt]
\textit{Ghi chú: Đồ thị $A$ và bản đánh số lại $PAP^{\top}$ thực chất là cùng một đồ thị, nên mô hình dự đoán ở mức đồ thị phải cho cùng kết quả, tức $f(PAP^{\top})=f(A)$.}}
\label{fig:bat-bien-do-thi}
\end{figure}
Lập luận tương tự áp dụng cho point cloud: một vật thể được mô tả bởi một tập vector toạ độ, và thứ tự liệt kê các vector này không mang ý nghĩa hình học nên mô hình cần bất biến với hoán vị các điểm.
\end{example}
 
\subsubsection{Tính đẳng biến / tương biến}
\label{ssub:equivariance}
 
Bất biến phù hợp khi đầu ra là một đại lượng toàn cục. Tuy nhiên nhiều bài toán có đầu ra mang cấu trúc như toạ độ, lực, mặt nạ phân đoạn. Khi đó ta không muốn đầu ra bất biến, mà muốn nó biến thiên một cách tương ứng với đầu vào.
 
\begin{definition}[Đẳng biến]
\label{def:dang-bien}
Hàm $f:\mathcal{X}\to\mathcal{Y}$ được gọi là đẳng biến đối với phép biến đổi $T$ nếu tồn tại phép biến đổi tương ứng $T'$ trên không gian đầu ra sao cho~\cite{cohen2016group}
\begin{equation}
\label{eq:equivariance}
f\bigl(T(x)\bigr) = T'\bigl(f(x)\bigr), \qquad \forall x\in\mathcal{X}.
\end{equation}
Trong trường hợp đặc biệt thường gặp khi $\mathcal{X}$ và $\mathcal{Y}$ chịu
cùng một phép biến đổi $T'=T$, điều kiện trở thành $f\bigl(T(x)\bigr)=T\bigl(f(x)\bigr)$: biến đổi đầu vào rồi đưa qua mô hình cho cùng kết quả với đưa qua mô hình rồi biến đổi đầu ra.
\end{definition}
 
\begin{example}[Nhúng nút trên đồ thị: đẳng biến hoán vị]
\label{vd:node-equiv}
Xét bài toán học nhúng cho từng nút bằng mạng nơ-ron đồ thị hoặc Transformer: đầu vào là tập đặc trưng nút, đầu ra là tập biểu diễn tương ứng theo cùng thứ tự. Vì thứ tự liệt kê nút là tuỳ ý nhưng đồ thị không đổi, nếu ta hoán vị đầu vào thì đầu ra phải bị hoán vị y hệt:
\begin{equation}
f(PX) = P f(X), \qquad \forall P\in\mathcal{P}_n .
\end{equation}
\end{example}

\begin{figure}[H]\centering
\begin{tikzpicture}[font=\small,
  gnn/.style={draw=gmlblue,fill=gmlblue,text=white,rounded corners=2pt,
              minimum width=1.1cm,minimum height=0.65cm,font=\bfseries\footnotesize},
  parr/.style={-{Latex[length=2.2mm]},gmlorange,very thick}]
\draw[black!30,line width=0.7pt](.5,.95)--(0,0)--(1,0)--cycle;
\filldraw[c1!20,draw=c1,line width=0.8pt](.5,.95)circle(.22);
\filldraw[c2!20,draw=c2,line width=0.8pt](0,0)circle(.22);
\filldraw[c3!20,draw=c3,line width=0.8pt](1,0)circle(.22);
\node[c1!70!black,font=\footnotesize\bfseries]at(.5,.95){1};
\node[c2!70!black,font=\footnotesize\bfseries]at(0,0){2};
\node[c3!70!black,font=\footnotesize\bfseries]at(1,0){3};
\node[font=\footnotesize]at(.5,-.40){$X$};
 
\draw[parr](1.55,.45)--(2.65,.45)node[midway,above,font=\footnotesize]{$P$};
\begin{scope}[shift={(3.2,0)}]
\draw[black!30,line width=0.7pt](.5,.95)--(0,0)--(1,0)--cycle;
\filldraw[c2!20,draw=c2,line width=0.8pt](.5,.95)circle(.22);
\filldraw[c1!20,draw=c1,line width=0.8pt](0,0)circle(.22);
\filldraw[c3!20,draw=c3,line width=0.8pt](1,0)circle(.22);
\node[c2!70!black,font=\footnotesize\bfseries]at(.5,.95){1};
\node[c1!70!black,font=\footnotesize\bfseries]at(0,0){2};
\node[c3!70!black,font=\footnotesize\bfseries]at(1,0){3};
\node[font=\footnotesize]at(.5,-.40){$PX$};
\end{scope}
 
\node[gnn](fL)at(.5,-1.2){GNN};
\node[gnn](fR)at(3.7,-1.2){GNN};
\draw[-{Latex[length=1.6mm]}](.5,-.48)--(fL.north);
\draw[-{Latex[length=1.6mm]}](3.7,-.48)--(fR.north);
\foreach \i/\col/\lab in {0/c1/$h_1$,1/c2/$h_2$,2/c3/$h_3$}{
  \filldraw[\col!22,draw=\col,line width=0.6pt,rounded corners=1pt]
    (-.18+\i*.40,-2.18)rectangle(.22+\i*.40,-1.82);
  \node[\col!70!black,font=\tiny]at(.02+\i*.40,-2.00){\lab};}
\draw[-{Latex[length=1.6mm]}](fL.south)--(.5,-1.82);
\node[font=\tiny]at(.5,-2.38){$H=f(X)$};
\foreach \i/\col/\lab in {0/c2/$h_2$,1/c1/$h_1$,2/c3/$h_3$}{
  \filldraw[\col!22,draw=\col,line width=0.6pt,rounded corners=1pt]
    (3.02+\i*.40,-2.18)rectangle(3.42+\i*.40,-1.82);
  \node[\col!70!black,font=\tiny]at(3.22+\i*.40,-2.00){\lab};}
\draw[-{Latex[length=1.6mm]}](fR.south)--(3.7,-1.82);
\node[font=\tiny]at(3.7,-2.38){$PH=f(PX)$};
 
\draw[parr](1.45,-2.00)--(2.65,-2.00)node[midway,above,font=\footnotesize]{$P$};
\end{tikzpicture}
\caption[Đẳng biến hoán vị của GNN]{Đẳng biến hoán vị của GNN.\protect\\[2pt]
\textit{Ghi chú: Khi hoán vị thứ tự nút ở đầu vào ($PX$), nhúng đầu ra bị hoán vị theo đúng thứ tự đó ($PH$), tức $f(PX)=P f(X)$. Màu sắc theo dõi danh tính nút xuyên suốt.}}
\label{fig:equiv-node}
\end{figure}
 
\begin{example}[Dự đoán lực tác dụng lên phân tử: tính đẳng biến]
\label{vd:molecule-equivariance}
Xét bài toán dự đoán các vector lực tác dụng lên từng nguyên tử của một phân tử hóa học. Nếu ta xoay toàn bộ cấu trúc phân tử này một góc $-90^\circ$ trong không gian, thì các vector lực tương ứng được biểu diễn bằng các mũi tên màu đỏ cũng phải xoay theo đúng một góc $-90^\circ$ tương ứng. Tính chất này tuân theo hệ thức đẳng biến chuẩn xác: $f(T(x)) = T(f(x))$, giúp mô hình giữ được sự nhất quán vật lý bất kể định hướng của phân tử trong không gian.
\end{example}

\begin{figure}[H]\centering
\begin{tikzpicture}[font=\small,
  farr/.style={-{Latex[length=2.2mm]},red!65!black,line width=1.1pt},
  garr/.style={-{Latex[length=2.4mm]},gmlorange,very thick}]
\newcommand{\water}{%
  \draw[gmlblue!40,line width=1.1pt](0,0)--(-.72,-.72);
  \draw[gmlblue!40,line width=1.1pt](0,0)--(.72,-.72);
  \node[circle, fill=gmlblue!20, draw=gmlblue, line width=0.9pt, minimum size=0.52cm, inner sep=0pt, gmlblue!80!black, font=\tiny\bfseries] (O) at (0,0) {O};
  \node[circle, fill=gmlblue!8, draw=gmlblue!60, line width=0.7pt, minimum size=0.36cm, inner sep=0pt, gmlblue!65!black, font=\tiny] (H1) at (-.72,-.72) {H};
  \node[circle, fill=gmlblue!8, draw=gmlblue!60, line width=0.7pt, minimum size=0.36cm, inner sep=0pt, gmlblue!65!black, font=\tiny] (H2) at (.72,-.72) {H};
  \draw[farr](O)--(.44,.60);      
  \draw[farr](H1)--(-1.26,-.98); 
  \draw[farr](H2)--(1.24,-.44); 
}
 
\begin{scope}[shift={(1.2,1.65)}]
  \water
\end{scope}
\node[font=\footnotesize]at(1.2,.35){$x$};
\node[red!62!black,font=\scriptsize]at(1.2,.08){lực $= f(x)$};
 
\draw[garr](2.55,1.1)--(3.85,1.1)node[midway,above,font=\footnotesize]{$T$};
\node[gmlorange,font=\scriptsize]at(3.20,.77){quay $-90^{\circ}$};
 
\begin{scope}[shift={(5.6,1.65)},rotate=-90]
  \water
\end{scope}
\node[font=\footnotesize]at(5.6,.35){$T(x)$};
\node[red!62!black,font=\scriptsize,align=center]at(5.6,.08){lực $=T(f(x))=f(T(x))$};
 
\node[gmlblue!70!black,font=\scriptsize]at(3.4,-.45)
  {$f(T(x))=T(f(x))$};
\end{tikzpicture}
\caption{Đẳng biến $SE(3)$ trong dự đoán lực phân tử.}
\vspace{2pt}
\begin{center}
\small\textit{Ghi chú: Khi quay toàn bộ phân tử bằng phép biến đổi $T$, các vector lực dự đoán (mũi tên đỏ) cũng quay theo đúng phép biến đổi đó, tức $f(T(x))=T(f(x))$.}
\end{center}
\label{fig:equiv-se3}
\end{figure}

\begin{infobox}{Bất biến là trường hợp đặc biệt của đẳng biến}
Cho hàm đẳng biến thỏa mãn phương trình tổng quát:
\begin{equation}
\label{eq:equivariance-special}
f(T(x)) = T'(f(x)) \quad \forall x \in \mathcal{X}
\end{equation}
Trong đó $T'$ là phép biến đổi tương ứng tác động lên không gian đầu ra.

Nếu ta xét trường hợp đặc biệt khi $T' = \mathrm{Id}$ tức phép đồng nhất, nghĩa là phép biến đổi ở đầu ra là tầm thường và không làm thay đổi bất kỳ phần tử nào:
\begin{equation}
T'(y) = \mathrm{Id}(y) = y \quad \forall y \in \mathcal{Y}
\end{equation}

Khi đó, thay $y = f(x)$ vào biểu thức trên, vế phải của phương trình~\eqref{eq:equivariance-special} trở thành:
\begin{equation}
T'(f(x)) = \mathrm{Id}(f(x)) = f(x)
\end{equation}

Thế ngược lại vào phương trình~\eqref{eq:equivariance-special}, ta thu được ngay:
\begin{equation}
\label{eq:invariance-special}
f(T(x)) = f(x) \quad \forall x \in \mathcal{X}
\end{equation}
Đây chính là định nghĩa của tính bất biến. 


\textbf{Kết luận:} Bất biến chính là đẳng biến với phép biến đổi đồng nhất ở đầu ra.
\end{infobox}
 
Mối quan hệ giữa biến đổi đầu vào và biến đổi đầu ra được tóm tắt gọn gàng bằng
sơ đồ giao hoán ở Hình~\ref{fig:commute}. Tính đẳng bất biến tương đương với việc sơ đồ tương ứng giao hoán: mọi đường đi từ góc trên-trái đến góc dưới-phải (hoặc tới $\mathcal{Y}$) đều cho cùng kết
quả.
 
\begin{figure}[H]
\centering
\begin{tikzpicture}[font=\normalsize]

  \node at (0,2.6) {\textbf{(a) Đẳng biến}};
  \node (X1) at (0,1.6) {$\mathcal{X}$};
  \node (Y1) at (2.4,1.6) {$\mathcal{Y}$};
  \node (X2) at (0,0) {$\mathcal{X}$};
  \node (Y2) at (2.4,0) {$\mathcal{Y}$};
  \draw[-{Latex[length=2mm]}] (X1) -- node[above]{$f$} (Y1);
  \draw[-{Latex[length=2mm]}] (X2) -- node[below]{$f$} (Y2);
  \draw[-{Latex[length=2mm]}] (X1) -- node[left]{$T$} (X2);
  \draw[-{Latex[length=2mm]}] (Y1) -- node[right]{$T'$} (Y2);

  \node at (7.2,2.6) {\textbf{(b) Bất biến}};
  \node (A1) at (6.0,1.6) {$\mathcal{X}$};
  \node (B1) at (8.4,0.8) {$\mathcal{Y}$};
  \node (A2) at (6.0,0) {$\mathcal{X}$};
  \draw[-{Latex[length=2mm]}] (A1) -- node[above right]{$f$} (B1);
  \draw[-{Latex[length=2mm]}] (A2) -- node[below right]{$f$} (B1);
  \draw[-{Latex[length=2mm]}] (A1) -- node[left]{$T$} (A2);
\end{tikzpicture}
\caption[Sơ đồ giao hoán của bất biến và đẳng biến]{Sơ đồ giao hoán của bất biến và đẳng biến.\protect\\[2pt]
\textit{Ghi chú: (a) Đẳng biến: phép biến đổi $T$ trên đầu vào tương ứng với phép biến đổi $T'$ trên đầu ra. (b) Bất biến: đầu ra không đổi dù đầu vào bị biến đổi bởi $T$.}}
\label{fig:commute}
\end{figure}
 
\subsubsection{Vì sao tính đối xứng lại quan trọng?}
Việc mã hoá sẵn đối xứng vào mô hình mang lại ba lợi ích đã được chứng thực
rộng rãi trong thực nghiệm:
\begin{itemize}[leftmargin=1.4em, itemsep=2pt, topsep=2pt]
  \item \textbf{Giảm bậc tự do và tăng hiệu quả dữ liệu.} Nếu đã biết hàm cần
  học có sẵn đối xứng, ta không cần học lại đối xứng đó từ dữ liệu. Trong bài toán dự đoán thế năng liên nguyên tử, các mô hình bất biến/đẳng biến đạt cùng độ chính xác trong khi dùng ít hơn tới ba bậc độ lớn khoảng $10^{3}$ lần dữ liệu huấn luyện so với mô hình không đối xứng~\cite{batzner2022nequip}.
  \item \textbf{Tăng độ bền vững và tổng quát hoá ngoài phân phối.} Nếu
  mô hình tôn trọng phép biến đổi, ta biết chính xác nó sẽ ngoại suy thế nào khi dữ liệu bị xoay/tịnh tiến mà không bị ảnh hưởng bởi thay đổi nhỏ. Vô hình chung là hành vi có thể dự đoán được và ổn định hơn.
  \item \textbf{Nền tảng cho mô hình tạo sinh.} Hiểu cơ chế duy trì và phá vỡ
  đối xứng cung cấp khung lý thuyết để thiết kế mô hình sinh cấu trúc phức tạp như phân tử thuốc, protein.
\end{itemize}
 
\subsection{Cơ sở lý thuyết nhóm}
\label{subsec:group-theory}
 
Các phép biến đổi mà ta quan tâm như phép quay, tịnh tiến, hoán vị không phải
là tuỳ ý: chúng có thể kết hợp với nhau, có thể đảo ngược, và luôn chứa phép không làm gì cả. Cấu trúc đại số nắm bắt chính xác những tính chất này là nhóm.
 
\begin{definition}[Nhóm]
\label{def:nhom}
Một nhóm là một cặp $(G,\cdot)$ gồm một tập hợp khác rỗng $G$ và một phép toán hai ngôi
$\cdot:G\times G\to G$ gọi là phép hợp thành, thoả mãn bốn tính chất~\cite{bronstein2021geometric}:
\begin{enumerate}[label=\textbf{(G\arabic*)}, leftmargin=2.6em, itemsep=2pt, topsep=2pt]
  \item \textbf{Tính đóng:} với mọi $g,h\in G$ thì $g\cdot h\in G$.
  \item \textbf{Tính kết hợp:}
        $(g\cdot h)\cdot \ell = g\cdot(h\cdot \ell)$ với mọi $g,h,\ell\in G$.
  \item \textbf{Phần tử đơn vị:} tồn tại $e\in G$ sao cho
        $e\cdot g = g\cdot e = g$ với mọi $g\in G$.
  \item \textbf{Phần tử nghịch đảo:} với mỗi $g\in G$, tồn tại
        $g^{-1}\in G$ sao cho $g\cdot g^{-1}=g^{-1}\cdot g = e$.
\end{enumerate}
Để gọn, ta thường viết $gh$ thay cho $g\cdot h$.
\end{definition}
 
\begin{remark}
Định nghĩa nhóm không đòi hỏi tính giao hoán $gh=hg$. Nếu nhóm thoả thêm
tính chất này, ta gọi nó là nhóm giao hoán. Phép quay trong mặt phẳng là giao hoán, nhưng phép quay trong không gian ba chiều và phép
hoán vị với $n\ge 3$ thì không.
\end{remark}
 
\begin{example}[Một số nhóm tiêu biểu]
\label{vd:cac-nhom}
\leavevmode
\begin{itemize}[leftmargin=1.4em, itemsep=2pt, topsep=2pt]
  \item \textbf{Nhóm tịnh tiến $(\mathbb{R}^{d},+)$:} Đại diện cho phép dịch chuyển vật thể theo một hướng và khoảng cách cố định mà không làm thay đổi hình dáng hay góc độ. Phần tử đơn vị là đứng yên, nghịch đảo của tiến lên là lùi lại.
  \item \textbf{Nhóm xoay tuần hoàn $C_{n}$:} Đại diện cho $n$ phép quay quanh một tâm cố định. Ví dụ $C_{4}$ là quay vật thể theo từng góc $90^{\circ}$ cho đến khi trở về vị trí ban đầu.
  \item \textbf{Nhóm Euclid $E(d)$ (và $SE(d)$):} Đại diện cho các phép dời hình cứng, kết hợp đồng thời cả việc dịch chuyển và xoay vật thể. Đây là nền tảng cốt lõi trong robotics và mô phỏng phân tử 3D.
\end{itemize}
\end{example}
 
\begin{figure}[H]
\centering
\begin{tikzpicture}[>=stealth, font=\small]
  \def\drawL#1#2#3{
    \draw[fill=#1, draw=#2, #3, thick, line join=round] 
      (0,0) -- (0.6,0) -- (0.6,0.2) -- (0.2,0.2) -- (0.2,0.8) -- (0,0.8) -- cycle;
  }
  \begin{scope}[xshift=-5cm]
    \begin{scope}[shift={(0,0)}]
      \drawL{cyan!40}{cyan!90}{}
    \end{scope}
  
    \begin{scope}[shift={(1.5, 1.2)}]
      \drawL{cyan!10}{cyan!50}{dashed}
    \end{scope}
    \draw[->, thick, red!70] (0.4, 0.4) -- (1.9, 1.6) node[midway, above left, text=black] {$v$};
    
    \node[below] at (1, -1.2) {\textbf{a)} Nhóm tịnh tiến};
  \end{scope}

  \begin{scope}[xshift=0cm, yshift=0.5cm]
    \fill[black!50] (0,0) circle (1.2pt) node[below right, inner sep=2pt] {tâm};
    \draw[black!20, dashed] (0,0) circle (1.2);
    \begin{scope}[shift={(1.2,0)}, rotate=0]     \begin{scope}[shift={(-0.3,-0.4)}] \drawL{orange!40}{orange!90}{} \end{scope} \end{scope}
    \begin{scope}[shift={(0,1.2)}, rotate=90]    \begin{scope}[shift={(-0.3,-0.4)}] \drawL{orange!10}{orange!50}{dashed} \end{scope} \end{scope}
    \begin{scope}[shift={(-1.2,0)}, rotate=180]  \begin{scope}[shift={(-0.3,-0.4)}] \drawL{orange!10}{orange!50}{dashed} \end{scope} \end{scope}
    \begin{scope}[shift={(0,-1.2)}, rotate=270]  \begin{scope}[shift={(-0.3,-0.4)}] \drawL{orange!10}{orange!50}{dashed} \end{scope} \end{scope}
    \draw[->, thick, red!70] (20:1.5) arc (20:70:1.5) node[midway, right=2pt, text=black] {$90^\circ$};
    
    \node[below] at (0, -1.7) {\textbf{b)} Nhóm xoay $C_4$};
  \end{scope}
  \begin{scope}[xshift=5cm]
    \begin{scope}[shift={(0,0)}]
      \drawL{teal!40}{teal!90}{}
    \end{scope}
    \begin{scope}[shift={(1.8, 1.1)}, rotate=-60]
      \drawL{teal!10}{teal!50}{dashed}
    \end{scope}
    \draw[->, thick, red!70] (0.4, 0.3) to[out=15, in=150] (1.6, 1.3) node[midway, below right, text=black] {$g \in SE(2)$};
    
    \node[below] at (1.2, -1.2) {\textbf{c)} Nhóm Euclid $SE(2)$};
  \end{scope}

\end{tikzpicture}
\caption[Minh hoạ tác động của ba nhóm tiêu biểu]{Minh hoạ tác động của ba nhóm tiêu biểu.\protect\\[2pt]
\textit{Ghi chú: (a) Nhóm tịnh tiến chỉ thay đổi vị trí; (b) nhóm xoay $C_4$ chỉ thay đổi góc xoay quanh tâm; (c) nhóm Euclid $SE(2)$ thay đổi đồng thời vị trí và góc xoay.}}
\label{fig:three_groups}
\end{figure}
 
\subsubsection{Tác động nhóm và biểu diễn nhóm}
\label{ssub:group-action}
 
Một nhóm trở nên hữu ích khi nó tác động lên dữ liệu: các phần tử của nhóm là những phép biến đổi, nên chúng phải biến đổi một thứ gì đó.
 
\begin{definition}[Tác động nhóm]
\label{def:tac-dong}
Một tác động của nhóm $G$ lên tập $\mathcal{X}$ là một ánh xạ $G\times\mathcal{X}\to\mathcal{X}$, $(g,x)\mapsto g\cdot x$, thoả mãn~\cite{bronstein2021geometric}:
\begin{equation}
e\cdot x = x \quad\text{và}\quad (gh)\cdot x = g\cdot(h\cdot x),
\qquad \forall g,h\in G,\ \forall x\in\mathcal{X}.
\end{equation}
Tức là phần tử đơn vị không làm gì, và tác động liên tiếp tương thích với phép
hợp thành trong nhóm.
\end{definition}
 
\begin{example}
\label{vd:action}
Hoán vị $\sigma\in S_{n}$ tác động lên tập chỉ số $\{1,\dots,n\}$ theo
$\sigma\cdot i = \sigma(i)$. Ma trận khả nghịch $A\in GL(n)$ tác động lên
$\mathbb{R}^{n}$ bằng phép nhân ma trận $A\cdot x = Ax$.
\end{example}
 
Trong các trường hợp ta quan tâm, tác động nhóm có thể được tuyến tính hoá: mỗi phần tử nhóm được biểu diễn bằng một ma trận. Đây là chiếc cầu nối giữa lý thuyết nhóm trừu tượng và đại số tuyến tính của mạng nơ-ron.
 
\begin{definition}[Biểu diễn nhóm]
\label{def:bieu-dien}
Một biểu diễn của nhóm $G$ trên không gian vector $V$ là một ánh xạ $\rho:G\to GL(V)$ tương thích với phép hợp thành~\cite{bronstein2021geometric}:
\begin{equation}
\label{eq:rep-hom}
\rho(gh) = \rho(g)\,\rho(h), \qquad \forall g,h\in G,
\end{equation}
khi đó tác động được viết dưới dạng nhân ma trận $g\cdot x = \rho(g)\,x$.
\end{definition}
 
\begin{example}
\label{vd:rep}
Hoán vị $\sigma$ tác động lên vector qua ma trận hoán vị $\sigma\cdot x = P_{\sigma}\,x$. Phép quay tác động qua ma trận quay $g\cdot x = R\,x = \rho(g)\,x$. Hình~\ref{fig:permutation} minh hoạ tác động của một hoán vị lên một mảng năm phần tử.
\end{example}
 
\begin{figure}[H]
\centering
\begin{tikzpicture}[font=\small]
  \def\sq{0.7}
  \foreach \i/\c in {0/blue!70, 1/teal!70, 2/green!60!black, 3/orange!80, 4/red!70}{
    \fill[\c] (\i*\sq, 1.2) rectangle (\i*\sq+\sq, 1.2+\sq);
    \draw[black!50] (\i*\sq, 1.2) rectangle (\i*\sq+\sq, 1.2+\sq);
  }
  \foreach \i/\c in {0/blue!70, 1/orange!80, 2/green!60!black, 3/red!70, 4/teal!70}{
    \fill[\c] (\i*\sq, 0) rectangle (\i*\sq+\sq, \sq);
    \draw[black!50] (\i*\sq, 0) rectangle (\i*\sq+\sq, \sq);
  }
  \foreach \a/\b in {0/0, 1/4, 2/2, 3/1, 4/3}{
    \draw[-{Latex[length=1.6mm]}, red!55, thick]
      (\a*\sq+\sq/2, 1.2) -- (\b*\sq+\sq/2, \sq);
  }
  \node[anchor=west] at (5*\sq+0.3, 1.55) {$x$};
  \node[anchor=west] at (5*\sq+0.3, 0.35) {$\sigma\cdot x = P_{\sigma}\,x$};
\end{tikzpicture}
\caption[Tác động hoán vị lên một mảng]{Tác động hoán vị lên một mảng.\protect\\[2pt]
\textit{Ghi chú: Hoán vị $\sigma\in S_{5}$ tráo đổi thứ tự các phần tử của mảng đầu vào $x$, biểu diễn bằng phép nhân với ma trận hoán vị $P_{\sigma}$.}}
\label{fig:permutation}
\end{figure}
 
\subsubsection{Quỹ đạo}
\label{ssub:orbit}
 
\begin{definition}[Quỹ đạo]
\label{def:quy-dao}
Quỹ đạo của một điểm $x\in\mathcal{X}$ dưới tác động của nhóm $G$ là tập hợp tất cả các điểm có thể đạt tới từ $x$ bằng các phép biến đổi trong $G$~\cite{jegelka2025tutorial}:
\begin{equation}
\label{eq:orbit}
\mathrm{Orb}_{G}(x) = \{\, g\cdot x \mid g\in G \,\}.
\end{equation}
\end{definition}
 
\begin{example}
\label{vd:orbit}
Với nhóm $C_{4}$ các phép quay $90^{\circ}$ tác động lên một điểm trong mặt phẳng, quỹ đạo gồm đúng bốn điểm là bốn vị trí quay của điểm đó (Hình~\ref{fig:orbit}). Với nhóm tịnh tiến tác động lên một ảnh, quỹ đạo của ảnh
là tập tất cả các phiên bản dịch chuyển của nó.
\end{example}
 
\begin{figure}[H]
\centering
\begin{tikzpicture}[font=\small]
  \def\rad{1.5}
  \draw[black!25, dashed] (0,0) circle (\rad);
  \fill[black!35] (0,0) circle (1.2pt);
  \node[anchor=north, font=\small] at (0, -0.15) {tâm quay};
  \foreach \ang/\lab/\col in {45/{$x$}/gmlblue, 135/{$g\cdot x$}/gmlorange,
                               225/{$g^{2}\!\cdot x$}/gmlorange, 315/{$g^{3}\!\cdot x$}/gmlorange}{
    \fill[\col] (\ang:\rad) circle (2.6pt);
    \node[\col!60!black] at (\ang:\rad+0.55) {\lab};
  }
  \foreach \a/\b in {45/135, 135/225, 225/315, 315/405}{
    \draw[-{Latex[length=1.8mm]}, red!55, thick] (\a+8:\rad) arc (\a+8:\b-8:\rad);
  }
  \node[red!60!black] at (90:\rad+0.3) {\scriptsize $g=90^{\circ}$};
  \node[anchor=west, align=left] at (2.4,0)
    {$\mathrm{Orb}_{C_{4}}(x)=$\\[2pt]$\{x,\,g\cdot x,\,g^{2}\!\cdot x,\,g^{3}\!\cdot x\}$};
\end{tikzpicture}
\caption[Quỹ đạo của một điểm dưới nhóm $C_4$]{Quỹ đạo của một điểm dưới nhóm $C_4$.\protect\\[2pt]
\textit{Ghi chú: Dưới tác động của nhóm xoay tuần hoàn $C_4$ (phép quay $90^{\circ}$), quỹ đạo của điểm $x$ gồm bốn vị trí quay liên tiếp.}}
\label{fig:orbit}
\end{figure}
 
\begin{infobox}{Liên hệ khái niệm quỹ đạo với tính bất biến}
Một hàm $f$ là bất biến $G$ khi và chỉ khi nó nhận giá trị hằng trên
mỗi quỹ đạo: $f$ không phân biệt được hai điểm thuộc cùng một
$\mathrm{Orb}_{G}(x)$. Quan sát đơn giản này là chìa khoá cho hai họ phương pháp
ở chương~3: canonicalization chọn một đại diện cố định cho mỗi quỹ đạo, hay group/frame averaging trung bình hoá giá trị trên toàn bộ quỹ đạo để ép buộc tính bất biến.
\end{infobox}

\subsubsection{Phát biểu hình thức của bất biến và đẳng biến}
\label{ssub:formal-inv-equiv}
 
Với ngôn ngữ tác động nhóm, hai khái niệm ở mục~\ref{subsec:doi-xung-bat-bien}
được phát biểu lại một cách chính xác và tổng quát.
 
\begin{definition}[Bất biến và đẳng biến theo nhóm]
\label{def:formal}
Cho nhóm $G$ tác động lên không gian đầu vào $\mathcal{X}$ và đầu ra
$\mathcal{Y}$. Hàm $f:\mathcal{X}\to\mathcal{Y}$ được gọi là:
\begin{itemize}[leftmargin=1.4em, itemsep=2pt, topsep=2pt]
  \item G-invariant nếu
        $\displaystyle f(g\cdot x) = f(x)$ với mọi $g\in G$;
  \item G-equivariant nếu
        $\displaystyle f(g\cdot x) = g\cdot f(x)$ với mọi $g\in G$, hay tổng quát
        hơn $f\bigl(\rho_{1}(g)\,x\bigr)=\rho_{2}(g)\,f(x)$ khi tác động trên
        đầu vào và đầu ra dùng hai biểu diễn $\rho_{1},\rho_{2}$ khác nhau.
\end{itemize}
\end{definition}
 
 
\begin{table}[H]
\centering
\caption{So sánh hai loại đối xứng cơ bản.}
\label{tab:inv-vs-equiv}
\renewcommand{\arraystretch}{1.3}
\begin{tabularx}{\linewidth}{@{}l L L@{}}
\toprule
\rowcolor{black!8} & \textbf{Bất biến } & \textbf{Đẳng biến}\\
\midrule
Điều kiện      & $f(g\cdot x)=f(x)$        & $f(g\cdot x)=g\cdot f(x)$ \\
Đầu ra điển hình & Vô hướng & Vector / có cấu trúc \\
Ví dụ          & Phân loại ảnh/đồ thị      & Nhúng nút, dự đoán lực, xác định khung bao \\
Quan hệ        & \multicolumn{2}{c}{Bất biến là trường hợp riêng của đẳng biến} \\
\bottomrule
\end{tabularx}
\end{table}

\begin{figure}[H]
\centering
\begin{tikzpicture}[font=\small]
  \node[anchor=west] at (-6.1,2.35)
    {\textbf{(a) Bất biến} (phân loại):\; $f(g\cdot x)=f(x)$};
  \draw[rounded corners=3pt, thick, black!55, fill=blue!3] (-6.0,0.5) rectangle (-4.6,1.9);
  \catface{-5.55}{0.90}{0.8}
  \draw[-{Latex[length=1.8mm]}, thick] (-4.52,1.2) -- (-3.78,1.2);
  \node at (-4.15,1.46) {$g$};
  \draw[rounded corners=3pt, thick, black!55, fill=blue!3] (-3.7,0.5) rectangle (-2.3,1.9);
  \catface{-2.70}{1.30}{0.8}
  \node[draw, rounded corners=2pt, fill=gmlblue, text=white,
        minimum width=0.8cm, minimum height=0.62cm] (fa) at (-1.0,1.2) {$f$};
  \draw[-{Latex[length=1.8mm]}, thick] (-2.22,1.2) -- (fa.west);
  \node[draw, rounded corners=2pt, fill=green!12, minimum width=1.0cm] (ya) at (1.1,1.2) {Mèo};
  \draw[-{Latex[length=1.8mm]}, thick] (fa.east) -- (ya.west);
  \node[anchor=west] at (-6.1,-0.25)
    {\textbf{(b) Đẳng biến} (định vị):\; $f(g\cdot x)=g\cdot f(x)$};
  \draw[rounded corners=3pt, thick, black!55, fill=blue!3] (-6.0,-1.9) rectangle (-4.6,-0.5);
  \catface{-5.60}{-1.20}{0.8}
  \draw[dashed, red!75, thick] (-5.95,-1.62) rectangle (-5.26,-0.78);
  \draw[-{Latex[length=1.8mm]}, thick] (-4.52,-1.2) -- (-3.78,-1.2);
  \node at (-4.15,-0.94) {$g$};
  \draw[rounded corners=3pt, thick, black!55, fill=blue!3] (-3.7,-1.9) rectangle (-2.3,-0.5);
  \catface{-2.66}{-1.20}{0.8}
  \draw[dashed, red!75, thick] (-3.01,-1.62) rectangle (-2.32,-0.78);
  \node[draw, rounded corners=2pt, fill=gmlblue, text=white,
        minimum width=0.8cm, minimum height=0.62cm] (fb) at (-1.0,-1.2) {$f$};
  \draw[-{Latex[length=1.8mm]}, thick] (-2.22,-1.2) -- (fb.west);
  \node[draw, rounded corners=2pt, fill=red!8, minimum width=1.5cm] (yb) at (1.35,-1.2)
        {hộp $g\cdot\mathrm{box}$};
  \draw[-{Latex[length=1.8mm]}, thick] (fb.east) -- (yb.west);
\end{tikzpicture}
\caption[Trực quan bất biến và đẳng biến]{Trực quan hoá bất biến và đẳng biến.\protect\\[2pt]
\textit{Ghi chú: (a) Bất biến: tác động $g$ lên con mèo trong khung hình không làm thay đổi nhãn đầu ra, $f(g\cdot x)=f(x)$. (b) Đẳng biến: khung bao đầu ra biến đổi tương ứng, $f(g\cdot x)=g\cdot f(x)$.}}
\label{fig:inv-equiv}
\end{figure}

\clearpage